\section{Introduction}
\label{sec:introduction}
Financial crime surveillance represents a vital, high-stakes application of intelligent expert systems within global financial infrastructure. According to the United Nations Office on Drugs and Crime \citep{unodc2011illicit}, illicit syndicates launder between \$800 billion and \$2 trillion annually (representing 2\% to 5\% of global Gross Domestic Product) across inter-bank wire rails, mobile financial service (MFS) e-wallets, and distributed crypto-asset ledgers. Operational Financial Intelligence Units (FIUs) within commercial banking institutions traditionally rely on rule-based alert systems. However, these legacy architectures generate excessive false positive alert rates exceeding 90\% to 95\%, imposing severe investigative backlogs and substantial operational expenses on compliance examiners.

To modernize financial surveillance, recent research has explored Graph Neural Networks (GNNs) to model the relational topologies of transaction graphs \citep{kipf2017semi, hamilton2017inductive, weber2019anti, hu2020heterogeneous}. Nevertheless, deploying automated graph intelligence in mission-critical Anti-Money Laundering (AML) production pipelines reveals five severe operational challenges:
\begin{enumerate}
    \item \textbf{Continuous-Time Velocity Evasion:} Laundering networks alternate between sub-second structuring bursts and multi-week dormant holding intervals, dynamics that cannot be captured simultaneously by discrete snapshot windowing \citep{pareja2020evolvegcn} or single-rate exponential decay kernels.
    \item \textbf{Adversarial Topological Camouflage:} Illicit entities deliberately route high-degree synthetic transaction chaff to legitimate commercial hubs (e.g., merchant gateways and centralized exchanges), inducing severe feature over-smoothing during neighborhood message aggregation \citep{dou2020enhancing}.
    \item \textbf{Extreme Class Imbalance ($\pi < 0.05\%$):} Extreme minority class skew induces severe posterior score attenuation, triggering near-complete recall collapse under standard decision cutoffs ($\tau = 0.50$).
    \item \textbf{Cold-Start Graph Isolation ($d(u) \le 2$):} Newly activated accounts possess sparse connectivity, receiving uninformative or noisy neighborhood aggregation vectors under pure spatial message passing.
    \item \textbf{Uncalibrated Decision Uncertainty:} Black-box anomaly scores lack distribution-free statistical risk guarantees under non-stationary temporal distribution drift, exposing financial institutions to regulatory non-compliance.
\end{enumerate}

To overcome these foundational hurdles, this study investigates three central research hypotheses:
\begin{itemize}
    \item ($H_1$) \textit{Camouflage Suppression:} Multi-scale continuous harmonic temporal attention coupled with learnable edge-trust filtering suppresses camouflage chaff across burst and dormancy cycles, mitigating feature over-smoothing and preserving minority discriminability relative to static and single-decay dynamic GNNs.
    \item ($H_2$) \textit{Flow-Conserving Latent Oversampling:} Interpolating synthetic minority instances strictly within latent topological typology clusters ($\mathcal{C}_k$) restores minority recall without elevating false alarms on commercial hubs, while formally preserving point-level fund conservation regularity.
    \item ($H_3$) \textit{Risk Guarantees:} Class-conditional Conformal Risk Control (CRC) establishes finite-sample label coverage ($\ge 99.0\%$) under within-class exchangeability with bounded temporal drift degradation, routing $>99.4\%$ of transaction volume into decisive singleton prediction sets.
\end{itemize}

A critical empirical finding of this work is that graph neural networks do not uniformly surpass optimized gradient-boosted decision trees across all financial transaction streams. Rather, their performance advantage is heavily concentrated in graph-rich, multi-hop relational laundering environments, whereas tree ensembles remain highly competitive or superior on predominantly flat, single-hop transactional ledgers.

Grounded in this operational reality, we propose \textbf{C-STGB} (\underline{C}onformal \underline{S}patio-\underline{T}emporal \underline{G}raph\underline{B}oost), an end-to-end intelligent surveillance and decision-support system. The principal scientific and systems contributions of this paper are fourfold:
\begin{enumerate}
    \item \textbf{Continuous-Time \& Camouflage-Resilient Architecture:} We develop an integrated spatio-temporal graph architecture combining continuous harmonic temporal encodings ($\Phi_{\text{time}}$) with Hawkes point-process arrival velocity descriptors and a context-aware edge-trust gate ($\hat{g}_{ij}$) that filters synthetic camouflage chaff and mitigates feature over-smoothing around high-degree banking hubs.
    \item \textbf{Flow-Conserving Typology Learning:} We design a typology-clustered latent Graph\-SMOTE mechanism paired with Asymmetric Focal Tversky optimization. We formally prove that synthetic node interpolation preserves point-level flow-conservation invariants while boosting minority detection under extreme class imbalance ($\pi < 0.05\%$).
    \item \textbf{Three-Tier Conformal Risk Guarantees:} We formulate Class-Conditional Conformal Risk Control for streaming transaction networks, providing finite-sample coverage ($\ge 99.0\%$) under exchangeability, non-asymptotic drift bounds under total variation shifts, and online Martingale tracking via Adaptive Conformal Inference (ACI).
    \item \textbf{Empirical Demarcation \& Production Systems Profiling:} Across 14 financial transaction networks comprising $9.53\text{M}$ entities, $9.53\text{M}$ transactions, and 182 evaluation pairs over five locked random seeds, we demonstrate statistically significant superiority over deep GNNs ($+43.49\text{ pp}$ mean uplift over GCN, $p_{\text{adj}} < 0.001$), define rigorous empirical boundaries where tabular learners remain competitive, and validate a streaming Fast-Path inference pipeline operating at $0.45$\,ms latency per single transaction.
\end{enumerate}

\textbf{Paper Organization:} The remainder of this paper is structured as follows: Section~\ref{sec:related_work} reviews related literature; Section~\ref{sec:methodology} formalizes the threat model and the C-STGB expert system architecture; Section~\ref{sec:experiments} presents the empirical evaluation across 14 benchmark networks; Section~\ref{sec:discussion} analyzes operational implications and domain boundaries; Section~\ref{sec:xai_case_study} details qualitative forensic attributions and compliance decision-support case studies; Section~\ref{sec:limitations} discusses limitations and ethical considerations; and Section~\ref{sec:conclusion} concludes the paper with future research directions.
